\section{Userspace ALSA}

\subsection{alsa-lib}

=== alsa-lib
  \begin{itemize}
  \item The main way to interact with ALSA devices is to use alsa-lib.
  \item \url{https://github.com/alsa-project/alsa-lib}
  \item It provides mainly access to the devices but also goes further
    and allows handling audio in userspace.
  \item The library itself is actually named ```libasound```
  \item The include file is ```alsa/asoundlib.h```
  \end{itemize}


=== [fragile]alsa-lib API
  \begin{itemize}
  \item 
    \begin{block}{}
    \fontsize{9}{9}\selectfont
      \begin{minted}{c}
int snd_pcm_open(snd_pcm_t ** pcmp, const char * name, snd_pcm_stream_t stream, int mode)
      \end{minted}
    \end{block}
  \item ```name``` is the name of the PCM to be opened.
  \item ```stream``` can be either ```SND_PCM_STREAM_PLAYBACK``` or
    ```SND_PCM_STREAM_CAPTURE```
  \item ```mode``` can be a combination of ```SND_PCM_NONBLOCK```
    and ```SND_PCM_ASYNC```
  \item 
    \begin{block}{}
    \fontsize{9}{9}\selectfont
      \begin{minted}{c}
int snd_pcm_close(snd_pcm_t *pcm)
      \end{minted}
    \end{block}
  \item Closes the PCM.
  \end{itemize}


=== PCM name
  \begin{itemize}
  \item This can be specified as a hardware device. The three
    arguments (in order: CARD,DEV,SUBDEV) specify card number or
    identifier, device number and subdevice number (-1 means any). For
    example: ```hw:0``` or ```hw:1,0```. Instead of the index, the
    card name can be used: ```hw:STM32MP15DK,0```
  \item Or through the ```plug``` plugin: ```plug:mypcmdef```,
    ```plug:hw:0,0```.
  \item The list of available names can be generated using
    ```snd_card_next``` to iterate over all the physical cards. See
    ```device_list``` in ```aplay```.
  \end{itemize}


=== [fragile]alsa-lib API - PCM
  \begin{itemize}
  \item The next step is to handle the PCM stream parameters
  \item
    \begin{block}{}
    \fontsize{9}{9}\selectfont
      \begin{minted}{c}
snd_pcm_hw_params_t *hw_params;
int snd_pcm_hw_params_malloc(snd_pcm_hw_params_t ** ptr)
int snd_pcm_hw_params_any(snd_pcm_t * pcm, snd_pcm_hw_params_t * params)
      \end{minted}
      \end{block}
    \item This will allocate a ```snd_pcm_hw_params_t``` and fill it
      with the range of parameters supported by ```pcm```.
  \item
    \begin{block}{}
    \fontsize{9}{9}\selectfont
      \begin{minted}{c}
int snd_pcm_hw_params_set_access(snd_pcm_t *pcm, snd_pcm_hw_params_t *params,
                                 snd_pcm_access_t _access)
      \end{minted}
    \end{block}
    \item This set the proper access type: interleaved or
      non-interleaved, mmap or not.
  \item
    \begin{block}{}
    \fontsize{9}{9}\selectfont
      \begin{minted}{c}
int snd_pcm_hw_params_set_format(snd_pcm_t *pcm, snd_pcm_hw_params_t *params,
                                 snd_pcm_format_t val)
      \end{minted}
    \end{block}
    \item This set the format, using a ```SND_PCM_FORMAT_``` macro.
  \end{itemize}


=== [fragile]alsa-lib API - PCM
  \begin{itemize}
  \item
    \begin{block}{}
    \fontsize{9}{9}\selectfont
      \begin{minted}{c}
int snd_pcm_hw_params_set_channels(snd_pcm_t *pcm, snd_pcm_hw_params_t *params,
                                   unsigned int val)
      \end{minted}
    \end{block}
    \item Sets the number of channels.
  \item
    \begin{block}{}
    \fontsize{9}{9}\selectfont
      \begin{minted}{c}
int snd_pcm_hw_params_set_rate_near(snd_pcm_t *pcm, snd_pcm_hw_params_t *params,
                                    unsigned int *val, int *dir)
      \end{minted}
    \end{block}
    \item Sets the sample rate, setting ```dir``` to 0 will require
      the exact rate.
  \item
    \begin{block}{}
    \fontsize{9}{9}\selectfont
      \begin{minted}{c}
int snd_pcm_hw_params_set_periods(snd_pcm_t *pcm, snd_pcm_hw_params_t *params,
                                  unsigned int val, int dir)
int snd_pcm_hw_params_set_period_size(snd_pcm_t *pcm, snd_pcm_hw_params_t *params,
                                      snd_pcm_uframes_t val, int dir)
int snd_pcm_hw_params_set_buffer_size(snd_pcm_t *pcm, snd_pcm_hw_params_t *params,
                                      snd_pcm_uframes_t val)
      \end{minted}
    \end{block}
  \item Sets the number of periods and the period size in the
      buffer or the buffer size.
  \end{itemize}


=== [fragile]alsa-lib API - PCM
  \begin{itemize}
  \item
    \begin{block}{}
    \fontsize{9}{9}\selectfont
      \begin{minted}{c}
int snd_pcm_hw_params(snd_pcm_t * pcm, snd_pcm_hw_params_t * params)
       \end{minted}
    \end{block}
  \item Installs the parameters and calls ```snd_pcm_prepare``` on
    the stream.
  \item
    \begin{block}{}
    \fontsize{9}{9}\selectfont
      \begin{minted}{c}
void snd_pcm_hw_params_free(snd_pcm_hw_params_t * obj)
      \end{minted}
    \end{block}
  \item Frees the allocated ```snd_pcm_hw_params_t```.
  \item
    \begin{block}{}
    \fontsize{9}{9}\selectfont
      \begin{minted}{c}
int snd_pcm_prepare(snd_pcm_t * pcm)
      \end{minted}
    \end{block}
  \item Prepares the stream.
  \item
    \begin{block}{}
    \fontsize{9}{9}\selectfont
      \begin{minted}{c}
int snd_pcm_wait(snd_pcm_t * pcm, int timeout)
      \end{minted}
    \end{block}
  \item Waits for the PCM to be ready.
  \end{itemize}


=== [fragile]alsa-lib API - PCM
  \begin{itemize}
  \item
    \begin{block}{}
    \fontsize{9}{9}\selectfont
      \begin{minted}{c}
snd_pcm_sframes_t snd_pcm_writei(snd_pcm_t *pcm, const void *buffer, snd_pcm_uframes_t size)
snd_pcm_sframes_t snd_pcm_readi(snd_pcm_t *pcm, void *buffer, snd_pcm_uframes_t size)
snd_pcm_sframes_t snd_pcm_writen(snd_pcm_t *pcm, void **bufs, snd_pcm_uframes_t size)
snd_pcm_sframes_t snd_pcm_readn(snd_pcm_t *pcm, void **bufs, snd_pcm_uframes_t size)
      \end{minted}
    \end{block}
  \item Write or read from an interleaved or non-interleaved buffer.
  \item
    \begin{block}{}
    \fontsize{9}{9}\selectfont
      \begin{minted}{c}
int snd_pcm_mmap_begin(snd_pcm_t *pcm, const snd_pcm_channel_area_t **areas,
                       snd_pcm_uframes_t *offset, snd_pcm_uframes_t *frames)
snd_pcm_sframes_t snd_pcm_mmap_commit(snd_pcm_t *pcm, snd_pcm_uframes_t offset,
                                      snd_pcm_uframes_t frames)
snd_pcm_sframes_t snd_pcm_mmap_writei(snd_pcm_t *pcm, const void *buffer,
                                      snd_pcm_uframes_t size)
snd_pcm_sframes_t snd_pcm_mmap_readi(snd_pcm_t *pcm, void *buffer, snd_pcm_uframes_t size)
snd_pcm_sframes_t snd_pcm_mmap_writen(snd_pcm_t *pcm, void **bufs, snd_pcm_uframes_t size)
snd_pcm_sframes_t snd_pcm_mmap_readn(snd_pcm_t *pcm, void **bufs, snd_pcm_uframes_t size)
      \end{minted}
    \end{block}
  \item Write or read from an interleaved or non-interleaved mmap buffer.
  \end{itemize}


=== [fragile]alsa-lib API - controls
  \begin{itemize}
  \item It is possible to set controls programatically.
  \item
    \begin{block}{}
    \fontsize{9}{9}\selectfont
      \begin{minted}{c}
snd_ctl_t *handle;
int snd_ctl_open (snd_ctl_t **ctl, const char *name, int mode)
      \end{minted}
    \end{block}
    \item Opens the sound card to be controlled.
  \item
    \begin{block}{}
    \fontsize{9}{9}\selectfont
      \begin{minted}{c}
snd_ctl_elem_id_t *id;
#define snd_ctl_elem_id_alloca(ptr)
snd_ctl_elem_value_t *value;
#define snd_ctl_elem_value_alloca(ptr)
      \end{minted}
    \end{block}
  \item Allocate a ```snd_ctl_elem_id_t```, referring to a particular
    control and a ```snd_ctl_elem_value_t``` to be set for this
    control.
  \item
    \begin{block}{}
    \fontsize{9}{9}\selectfont
      \begin{minted}{c}
void snd_ctl_elem_id_set_interface(snd_ctl_elem_id_t *obj, snd_ctl_elem_iface_t val)
void snd_ctl_elem_id_set_name(snd_ctl_elem_id_t *obj, const char *val)
      \end{minted}
    \end{block}
    \item Set the interface and name of the control to be set.
  \end{itemize}


=== [fragile]alsa-lib API - controls
  \begin{itemize}
  \item A lookup is needed to fill the ```snd_ctl_elem_id_t```
    completely
  \end{itemize}
  \begin{block}{}
    \fontsize{9}{9}\selectfont
    \begin{minted}{c}
int lookup_id(snd_ctl_elem_id_t *id, snd_ctl_t *handle)
{
    int err;
    snd_ctl_elem_info_t *info;
    snd_ctl_elem_info_alloca(&info);

    snd_ctl_elem_info_set_id(info, id);
    if ((err = snd_ctl_elem_info(handle, info)) < 0) {
        return err;
    ```
    snd_ctl_elem_info_get_id(info, id);

    return 0;
```
    \end{minted}
  \end{block}


=== [fragile]alsa-lib API - controls
  \begin{itemize}
  \item
    \begin{block}{}
    \fontsize{9}{9}\selectfont
      \begin{minted}{c}
void snd_ctl_elem_value_set_id(snd_ctl_elem_value_t *obj, const snd_ctl_elem_id_t *ptr)
      \end{minted}
    \end{block}
    \item Links the value with the control id.
  \item
    \begin{block}{}
    \fontsize{9}{9}\selectfont
      \begin{minted}{c}
void snd_ctl_elem_value_set_boolean(snd_ctl_elem_value_t *obj, unsigned int idx, long val)
void snd_ctl_elem_value_set_integer(snd_ctl_elem_value_t *obj, unsigned int idx, long val)
void snd_ctl_elem_value_set_integer64(snd_ctl_elem_value_t *obj, unsigned int idx,
                                      long long val)
void snd_ctl_elem_value_set_enumerated(snd_ctl_elem_value_t *obj, unsigned int idx,
                                       unsigned int val)
void snd_ctl_elem_value_set_byte(snd_ctl_elem_value_t *obj, unsigned int idx,
                                 unsigned char val)
void snd_ctl_elem_set_bytes(snd_ctl_elem_value_t *obj, void *data, size_t size)
      \end{minted}
    \end{block}
  \item Set the value in ```snd_ctl_elem_value_t```.
  \item
    \begin{block}{}
    \fontsize{9}{9}\selectfont
      \begin{minted}{c}
int snd_ctl_elem_write(snd_ctl_t *ctl, snd_ctl_elem_value_t *data)
      \end{minted}
    \end{block}
  \item Actually set the control.
  \end{itemize}


=== Going further
  \begin{itemize}
  \item UCM: The ALSA Use Case Configuration:
    \url{https://www.alsa-project.org/alsa-doc/alsa-lib/group__ucm__conf.html}
  \item ALSA topology:
    \url{https://www.alsa-project.org/wiki/ALSA_topology}
  \end{itemize}


\setupdemoframe{Card configuration examples}
```
  Using ```alsa-lib``` tools to:
  \begin{itemize}
  \item Reorder channels
  \item Split channels
  \item Resample
  \item Mix samples
  \item Apply effects
  \end{itemize}
```
